\chapter{Model, sectors, and bit-string bases}
\label{chap:basis}

\section{The reference problem}

For a chain of $L$ sites, the spin-$1/2$ Hubbard Hamiltonian is
\begin{equation}
H = -t\sum_{\langle i,j\rangle,\sigma}
\left(c^{\dagger}_{i\sigma}c_{j\sigma}
    +c^{\dagger}_{j\sigma}c_{i\sigma}\right)
+U\sum_{i=0}^{L-1}n_{i\uparrow}n_{i\downarrow}.
\label{eq:hubbard}
\end{equation}
The solver supports open and periodic nearest-neighbor chains.  Both the
up-spin number $N_\uparrow$ and down-spin number $N_\downarrow$ commute with
$H$, so the full Fock space splits into independent blocks.  We solve one block
at a time.

Working in a fixed sector is the first important reduction.  Instead of a
$4^L$-dimensional site Fock space, the sector dimension is
\begin{equation}
\basisdim(L,N_\uparrow,N_\downarrow)
= \binom{L}{N_\uparrow}\binom{L}{N_\downarrow}.
\label{eq:dimension}
\end{equation}
The reduction is exact because the Hamiltonian never changes either particle
number.

\section{Two integers per basis state}

A basis state is represented by
\begin{equation}
s=(b_\uparrow,b_\downarrow),
\end{equation}
where bit $i$ of $b_\sigma$ is one exactly when site $i$ contains a spin
$\sigma$ electron.  Site zero is the least significant bit.  For example, at
$L=4$,
\begin{equation}
(b_\uparrow,b_\downarrow)=(0101_2,0010_2)
\end{equation}
has up electrons on sites $0$ and $2$, and a down electron on site $1$.

This representation makes the common operations direct:
\begin{align}
n_{i\uparrow}(s) &= (b_\uparrow \mathbin{\texttt{>>}} i)\mathbin{\texttt{\&}}1,\\
N_\uparrow(s) &= \operatorname{popcount}(b_\uparrow),\\
N_{\mathrm{double}}(s) &=
\operatorname{popcount}(b_\uparrow\mathbin{\texttt{\&}}b_\downarrow).
\end{align}
Python's integer \code{bit\_count()} supplies the popcount operation.

\section{Enumeration and lookup}

All length-$L$ integers with popcount $N_\uparrow$ are combined with all
length-$L$ integers with popcount $N_\downarrow$.  A tuple stores the resulting
states in deterministic order, while a dictionary maps each state back to its
integer basis index.  The tuple gives $O(1)$ index-to-state lookup and the
dictionary gives average $O(1)$ state-to-index lookup.

The following first program verifies the combinatorial dimension, particle
counts, and round-trip lookup.

\lstinputlisting[caption={Enumerating a fixed sector.},label={lst:basis}]{code/step01_basis.py}

\section{Validation before allocation}

The basis constructor rejects $L<1$, particle numbers outside $[0,L]$, and
noninteger inputs.  It computes Eq.~\eqref{eq:dimension} before allocating the
state tuple.  A configurable dimension guard then refuses unreasonable
sectors.  This ordering matters: a friendly error is better than allowing the
operating system to terminate an overcommitted process.

At half filling with $N_\uparrow=N_\downarrow=L/2$, growth is already severe:
\begin{center}
\begin{tabular}{rrr}
\toprule
$L$ & $(N_\uparrow,N_\downarrow)$ & $\basisdim$ \\
\midrule
6  & $(3,3)$ & 400 \\
8  & $(4,4)$ & 4,900 \\
10 & $(5,5)$ & 63,504 \\
12 & $(6,6)$ & 853,776 \\
\bottomrule
\end{tabular}
\end{center}
The Hamiltonian need not be dense, but the basis and reverse-lookup dictionary
must still fit in memory.

\paragraph{Checkpoint.}
Before writing any operator code, verify for several sectors that every state
has the requested two popcounts, every state is unique, both lookup directions
round-trip, and the dimension equals Eq.~\eqref{eq:dimension}.
